%%%%%%%%%%%%%%%%%%%%%%%%%%%%%%%%%%%%%%%%%%%%%%%%%%%%%%%%%%%%%%%%%%%%%%%%%%%%
%% Author template for Organization Science (orsc)
%% Mirko Janc, Ph.D., INFORMS, mirko.janc@informs.org
%% ver. 0.95, December 2010
%% Adjusted by Nicolas Mangin (2018) to add a few environments
%%%%%%%%%%%%%%%%%%%%%%%%%%%%%%%%%%%%%%%%%%%%%%%%%%%%%%%%%%%%%%%%%%%%%%%%%%%%
%\documentclass[orsc,blindrev]{informs3}
\documentclass[$informsjnl$,$layout$]{fmt/informs3} % current default for manuscript submission

\DoubleSpacedXI % Made default on 4/4/2014 per request
%%\OneAndAHalfSpacedXI % current default line spacing
%%\OneAndAHalfSpacedXII
%%\DoubleSpacedXII

% If hyperref is used, dvi-to-ps driver of choice must be declared as
%   an additional option to the \documentclass. For example
%\documentclass[dvips,orsc]{informs3}      % if dvips is used
%\documentclass[dvipsone,orsc]{informs3}   % if dvipsone is used, etc.

%%% ORSC uses endnotes!
\usepackage{endnotes}
\let\footnote=\endnote
\let\enotesize=\normalsize
\def\notesname{Endnotes}%
\def\enoteformat{\rightskip0pt\leftskip0pt\parindent=0em
  \leavevmode\hbox{\makeenmark}}




% Additional libraries
\usepackage{adjustbox}
\usepackage{threeparttablex}
\usepackage{subcaption}
\usepackage{dcolumn}
\usepackage{tabularx}
\usepackage{longtable}
\usepackage{pdflscape}
\usepackage{booktabs}
\usepackage{geometry}
\usepackage[table]{xcolor}
\usepackage[figuresright]{rotating}


% Private macros here (check that there is no clash with the style)
% Solves the problem of tigthlists
\providecommand{\tightlist}{%
  \setlength{\itemsep}{0pt}\setlength{\parskip}{0pt}}

% Places all figures and charts at end of manuscript and adds 'insert table x about here' lines.
$if(tabfigtoend)$
  
  \usepackage[nolists]{endfloat} 
  \DeclareDelayedFloatFlavour*{sidewaystable}{table}
  \DeclareDelayedFloatFlavour*{longtable}{table}
  \DeclareDelayedFloatFlavour*{longsidewaystable}{table}
  \DeclareDelayedFloatFlavour*{landscape}{table}
  
  \renewcommand{\figureplace}{
    ------------------------------------\\
    Insert \figurename \ \thepostfig\ about here.\\
    ------------------------------------
    }
  \renewcommand{\tableplace}{
    ------------------------------------\\
    Insert \tablename \ \theposttbl\ about here.\\
    ------------------------------------
    }
$endif$

\usepackage[hidelinks]{hyperref}
  
% Defines environment for questions
\newtheorem{question}{Question}
\makeatletter
\newcounter{quest}
\newenvironment{gpquestion}
 {
  \refstepcounter{question}
  \protected@edef\theprop{\thequestion}
  \setcounter{quest}{\value{question}}
  \setcounter{question}{0}
  \def\thequestion{\theprop\alph{question}}
  \ignorespaces
}{
  \setcounter{question}{\value{quest}}
  \ignorespacesafterend
}
\makeatother

% Defines environment for hypotheses
\usepackage{amsmath}
\makeatletter
\newcounter{hyp}
\newenvironment{gphypothesis}
 {
  \refstepcounter{hypothesis}
  \protected@edef\thehyp{\thehypothesis}
  \setcounter{hyp}{\value{hypothesis}}
  \setcounter{hypothesis}{0}
  \def\thehypothesis{\thehyp\alph{hypothesis}}
  \ignorespaces
}{
  \setcounter{hypothesis}{\value{hyp}}
  \ignorespacesafterend
}
\makeatother

% Defines environment for propositions
\makeatletter
\newcounter{prop}
\newenvironment{gpproposition}
 {
  \refstepcounter{proposition}
  \protected@edef\theprop{\theproposition}
  \setcounter{prop}{\value{proposition}}
  \setcounter{proposition}{0}
  \def\theproposition{\theprop\alph{proposition}}
  \ignorespaces
}{
  \setcounter{proposition}{\value{prop}}
  \ignorespacesafterend
}
\makeatother

% Define the quote environment
\renewenvironment{quote}
               {\list{}{\rightmargin=3em\leftmargin=3em}
                \item\relax\itshape\ignorespaces}
               {\unskip\unskip\endlist}

\makeatletter
\newenvironment{quotep}[2][2em]
  {\noindent\setlength{\@tempdima}{#1}%
   \def\quotep@author{#2}%
   \parshape 1 \@tempdima \dimexpr\textwidth-2\@tempdima\relax\itshape\ignorespaces}
  {\par\normalfont\hfill--\ \quotep@author\hspace*{\@tempdima}\par\bigskip}
\makeatother

% Natbib setup for author-year style
\usepackage{natbib}
 \bibpunct[, ]{(}{)}{,}{a}{}{,}%
 \def\bibfont{\small}%
 \def\bibsep{\smallskipamount}%
 \def\bibhang{24pt}%
 \def\newblock{\ }%
 \def\BIBand{and}%

%% Setup of theorem styles. Outcomment only one. 
%% Preferred default is the first option.
\TheoremsNumberedThrough     % Preferred (Theorem 1, Lemma 1, Theorem 2)
%\TheoremsNumberedByChapter  % (Theorem 1.1, Lema 1.1, Theorem 1.2)

%% Setup of the equation numbering system. Outcomment only one.
%% Preferred default is the first option.
\EquationsNumberedThrough    % Default: (1), (2), ...
%\EquationsNumberedBySection % (1.1), (1.2), ...

% In the reviewing and copyediting stage enter the manuscript number.
\MANUSCRIPTNO{$submissionid$} % When the article is logged in and DOI assigned to it,
                 %   this manuscript number is no longer necessary

%%%%%%%%%%%%%%%%
\begin{document}
%%%%%%%%%%%%%%%%

% Outcomment only when entries are known. Otherwise leave as is and 
%   default values will be used.
%\setcounter{page}{1}
%\VOLUME{00}%
%\NO{0}%
%\MONTH{Xxxxx}% (month or a similar seasonal id)
%\YEAR{0000}% e.g., 2005
%\FIRSTPAGE{000}%
%\LASTPAGE{000}%
%\SHORTYEAR{00}% shortened year (two-digit)
%\ISSUE{0000} %
%\LONGFIRSTPAGE{0001} %
%\DOI{10.1287/xxxx.0000.0000}%

% Author's names for the running heads
% Sample depending on the number of authors;
% \RUNAUTHOR{Jones}
% \RUNAUTHOR{Jones and Wilson}
% \RUNAUTHOR{Jones, Miller, and Wilson}
% \RUNAUTHOR{Jones et al.} % for four or more authors
% Enter authors following the given pattern:
%\RUNAUTHOR{}

% Title or shortened title suitable for running heads. Sample:
% \RUNTITLE{Bundling Information Goods of Decreasing Value}
% Enter the (shortened) title:
%\RUNTITLE{}

% Full title. Sample:
\TITLE{$title$}
% Enter the full title:
%\TITLE{}

% Block of authors and their affiliations starts here:
% NOTE: Authors with same affiliation, if the order of authors allows, 
%   should be entered in ONE field, separated by a comma. 
%   \EMAIL field can be repeated if more than one author
\ARTICLEAUTHORS{%
$for(author)$
  \AUTHOR{$author.name$}
  \AFF{$author.affiliation$, $author.address$, \EMAIL{author.email}}
$endfor$
% Enter all authors
} % end of the block

\ABSTRACT{%
$abstract$ % Enter your abstract
}%

% Sample
%\KEYWORDS{deterministic inventory theory; infinite linear programming duality; 
%  existence of optimal policies; semi-Markov decision process; cyclic schedule}

% Fill in data. If unknown, outcomment the field
\KEYWORDS{$keywords$}

\maketitle
%%%%%%%%%%%%%%%%%%%%%%%%%%%%%%%%%%%%%%%%%%%%%%%%%%%%%%%%%%%%%%%%%%%%%%

% Samples of sectioning (and labeling) in ORSC
% NOTE: (1) \section and \subsection do NOT end with a period
%       (2) \subsubsection and lower need end punctuation
%       (3) capitalization is as shown (title style).
%
%\section{Introduction.}\label{intro} %%1.
%\subsection{Duality and the Classical EOQ Problem.}\label{class-EOQ} %% 1.1.
%\subsection{Outline.}\label{outline1} %% 1.2.
%\subsubsection{Cyclic Schedules for the General Deterministic SMDP.}
%  \label{cyclic-schedules} %% 1.2.1
%\section{Problem Description.}\label{problemdescription} %% 2.

% Text of your paper here

$body$

% Acknowledgments here
\ACKNOWLEDGMENT{%
$acknowledgements$ % Enter the text of acknowledgments here
}% Leave this (end of acknowledgment)


% Appendix here
% Options are (1) APPENDIX (with or without general title) or 
%             (2) APPENDICES (if it has more than one unrelated sections)
% Outcomment the appropriate case if necessary
%
% \begin{APPENDIX}{<Title of the Appendix>}
% \end{APPENDIX}
%
%   or 
%
% \begin{APPENDICES}
% \section{<Title of Section A>}
% \section{<Title of Section B>}
% etc
% \end{APPENDICES}

\bigskip

% Endnotes here
\theendnotes

\bigskip

% References here (outcomment the appropriate case) 

% CASE 1: BiBTeX used to constantly update the references 
%   (while the paper is being written).
\bibliographystyle{fmt/informs2014} % outcomment this and next line in Case 1
\bibliography{$bibliography$} % if more than one, comma separated

% CASE 2: BiBTeX used to generate mypaper.bbl (to be further fine tuned)
%\input{mypaper.bbl} % outcomment this line in Case 2


$for(include-after)$
$include-after$
$endfor$



\end{document}

